% GENERATED FILE. Edit canonical lessons, then run npm run generate:books.
% canonical-source-hash: fnv1a64:c7ec2578d7ec1f48
% canonical-lessons: ES-C53-que-relativo

\chapter{El Libro Que Compro --- One Sentence Out of Two}
\label{ch:relative-que}
\hlchaptermodality{182}

\begin{chapteropening}
\textbf{By the end of this chapter:} \emph{I can join two sentences with que so a noun carries a clause.}

Complete ES-C53-que-relativo: join two sentences with que and name the question words that share its root.
\end{chapteropening}

\section[el libro que compro]{el libro que compro — one sentence out of two}

\label{lesson:ES-C53-que-relativo}



\begin{quote}
Say \textquotedblleft{}I want you to speak.\textquotedblright{} (\emph{Quiero que hables}.) That \emph{que} is about to do a second job — and it is the job that lets your sentences stop being short.
\end{quote}

\begin{grammarlens}[title={Grammar Lens: stapling two sentences together}]
You have two sentences about one book:

\begin{quote}
\emph{Tengo el libro.} — I have the book. \emph{Compro el libro.} — I buy the book.
\end{quote}

\emph{Que} joins them, and the repeated noun disappears:

\begin{quote}
\emph{Tengo el libro \textbf{que} compro.} — I have the book \textbf{that} I buy.
\end{quote}

That is the whole mechanism. \textbf{Name the thing once, then use \emph{que} to hang a second sentence off it.} Everything you can already say can now be attached to a noun instead of standing alone.

This is the point where a language stops being a list of statements. Up to now your sentences have had one verb each. From here a sentence can contain another sentence, and then another — which is what allows description, explanation and storytelling of any length.
\end{grammarlens}

\begin{cousinweb}
\emph{Que} comes from the Latin \textbf{qui / quem / quid} family — the same \emph{qu-} question-root sitting inside \emph{qué}, \emph{quién}, \emph{cómo} and \emph{cuándo}, which you have been collecting since your very first questions.

That is not a coincidence of spelling. A question and a relative clause are doing related work: \emph{¿qué compras?} asks you to fill a gap, and \emph{el libro que compro} fills one. Latin used the same word family for both, and Spanish kept the arrangement.

So the busiest structural word in the language is a worn-down question word. English did the identical thing: \emph{what}, \emph{which}, \emph{who} and \emph{that} all serve twice, as questions and as joints.
\end{cousinweb}

\subsection*{Your turn}
Say these aloud:

\begin{itemize}
\raggedright
  \item the two short sentences, then the joined one
  \item \textquotedblleft{}el libro que compro\textquotedblright{}, \textquotedblleft{}la casa que tengo\textquotedblright{}
\end{itemize}

\begin{tcolorbox}[breakable,colback=teal!4,colframe=teal!35!black,title={Before you move on}]
Join \textquotedblleft{}I have the book\textquotedblright{} and \textquotedblleft{}I buy the book.\textquotedblright{} (\emph{\textbf{Tengo el libro que compro}}.) What does \emph{que} let a sentence do that it could not before? (\textbf{Contain another sentence.}) Next: the one place English is lazier than Spanish, and it trips every learner.
\end{tcolorbox}
